% ch2_theory.tex — 2장 이론적 기반과 해석적 맥락
% Called by: main.tex line 34 (% ch2_theory.tex — 2장 이론적 기반과 해석적 맥락
% Called by: main.tex line 34 (% ch2_theory.tex — 2장 이론적 기반과 해석적 맥락
% Called by: main.tex line 34 (% ch2_theory.tex — 2장 이론적 기반과 해석적 맥락
% Called by: main.tex line 34 (\input{sections/ch2_theory.tex})
\section{이론적 기반과 해석적 맥락}
\label{sec:theory}

\noindent{\sffamily\small\color{muted}\textbf{CHAPTER 2}}

\begin{quoteblock}
사람 기반과 장소 기반 접근을 통합하는 장소민감적(place-sensitive) 접근의 학술\textperiodcentered{}정책적 근거를 정리하고, 서울의 인구\textperiodcentered{}인프라\textperiodcentered{}정책 환경을 분석의 해석적 맥락으로 제시한다.
\end{quoteblock}

본 장은 분석에 관한 이론적 기반과 3장의 해석적 맥락을 제시한다. 먼저, 디지털 격차에 영향을 미치는 개인적 요인과 지역적 요인을 정리한다. 이후, 공공정책에서 사람 기반 접근(people-based)과 장소 기반 접근(place-based)의 구분과 최근 장소민감적 접근의 논의배경을 검토한다. 그 다음으로, 본 프로젝트의 메타이론적 틀인 비판적 실재론과 실제 연구 적용 사례를 검토한다. 마지막으로 서울의 인구, 행정, 인프라 조건을 검토해 분석 결과의 해석에 필요한 맥락 정보를 제공한다.

\subsection{Policy}
\label{subsec:2-1}

공공정책의 설계에서 개입의 대상을 사람의 특성에 두는 접근(people-based approach)과 장소의 조건에 두는 접근(place-based approach)의 구분은 오래된 논쟁이다. 사람 기반 정책은 개인의 역량, 소득, 교육 등 인적 특성을 직접 지원함으로써 결과를 개선하고자 한다. 장소 기반 정책은 특정 지역의 인프라, 제도, 환경 조건을 개선함으로써 그 지역 거주자의 결과를 간접적으로 향상시키고자 한다~\citep{NeumarkSimpson2015}. 두 접근은 상호 배타적이지 않으나, 정책 자원의 배분과 개입의 경로가 상이하므로 정책 수행의 관점으로 분리해 설계될 필요가 있다.

\subsubsection{디지털 격차의 영향 요인: 개인 특성과 지역 조건}
\label{subsubsec:2-1-1}

디지털 격차 문헌에서 개인 특성에 주목하는 연구는, 동일 지역 내에서도 개인 특성이 디지털 활용 격차를 설명하는 주된 요인임을 보여준다. \citet{Hargittai2010}는 인터넷 접근 이후의 격차, 즉 활용 방식의 차이가 실질적 불평등을 결정한다고 제시하였다. 사회경제적 지위에 따라 디지털 활용의 격차가 지속될 수 있음을 시사하는 결과이다. \citet{ChipevaEtAl2018}은 다양한 국가에서 실용성이 기술 사용의 가장 강력한 동기임을 보이는 한편, 개인의 내적 특성이 디지털 기기 사용에 유의미한 영향을 미친다는 점을 제시하였다. \citet{BorgesEtAl2020}은 앙골라 사례를 통해 기기 보급 이전에 습관과 ICT 기술 숙련도가 격차의 핵심 요인임을 보고하였다. \citet{Hilbert2011}의 라틴아메리카 18개국 분석은 남녀 간 디지털 격차가 존재하며, 국가 수준의 젠더 평등 지수가 높을수록 인터넷 사용이 활발함을 보였다.

한편, 개인 특성을 통제한 이후에도 지역 환경이 디지털 활용에 대해 추가적 설명력을 갖는다는 증거도 축적되어 있다. \citet{MossbergerEtAl2006}은 미국을 대상으로 HLM을 적용하여, 인터넷 접근 격차가 인종 자체보다 고빈곤 지역 환경에 기인함을 보였다. 이는 지역의 구조적 조건이 개인의 결과에 독립적 영향을 미칠 수 있음을 시사한다. \citet{CrangEtAl2006}은 영국 뉴캐슬의 질적 사례 비교를 통해, 지역 환경이 디지털 네트워크의 사용 방식 자체를 구조화함을 제시하였다. \citet{HongEtAl2017}의 중국 45세 이상 성인 분석은, 인터넷 이용자가 고소득, 고학력, 도시 거주 집단에 집중되어 있음을 보여, 지역 수준의 인프라와 기회 구조가 이용 가능성을 선별할 수 있음을 시사하였다.

\subsubsection{정책적 구분: 사람 기반 접근, 장소 기반 접근, 장소민감적 접근}
\label{subsubsec:2-1-2}

앞 절에서 검토한 두 유형의 요인은 공공정책 설계에서 사람 기반 접근(people-based approach)과 장소 기반 접근(place-based approach)의 구분에 대응한다. 사람 기반 정책은 개인의 역량, 소득, 교육 등 인적 특성을 직접 지원함으로써 결과를 개선하고자 한다. 장소 기반 정책은 특정 지역의 인프라, 제도, 환경 조건을 개선함으로써 그 지역 거주자의 결과를 간접적으로 향상시키고자 한다~\citep{NeumarkSimpson2015}. 두 접근은 상호 배타적이지 않으나, 정책 자원의 배분과 개입의 경로가 상이하므로 정책 수행의 관점에서 분리하여 설계될 필요가 있다.

그러나 두 접근 사이의 관계는 단순한 병렬이 아니다. \citet{Hindman2000}의 미국 시계열 분석은 소득, 연령, 교육이 지리보다 강력한 예측변수임을 보이면서도, 시간 경과에 따라 격차가 고착화되는 패턴을 보고하였다. 이는 개인 특성과 지역 구조가 시간 속에서 상호작용할 수 있음을 의미한다. \citet{vanDijk2020}는 동기 결여에서 물질적 접근, 기술 능력, 실질적 활용으로 이어지는 순차적 모형을 제시하여, 개인 요인과 구조 요인이 누적적으로 작용함을 제안하였다.

이러한 맥락에서, 최근 정책 및 학술 문헌은 사람 기반 또는 장소 기반의 양자택일이 아닌, 장소의 맥락 안에서 사람의 경험을 이해하는 장소민감적(place-sensitive) 접근으로 이행하고 있다. 생활 경험(lived experience)과 장소 기반 통찰(place-based insight)을 정량적 증거와 동등한 지식으로 인정하려는 시도가 그 예이다~\citep{OECD2025}. 이 접근은 사람 기반 요인과 장소 기반 요인을 분석적으로 분리하되, 두 요인이 특정 장소에서 어떻게 결합하는지를 함께 고려할 것을 요구한다.

\subsubsection{메타이론적 틀: 비판적 실재론의 층화된 존재론}
\label{subsubsec:2-1-3}

본 프로젝트는 사람 기반 요인과 장소 기반 요인의 분리를 분석 설계의 기본 원칙으로 채택하며, 그 설계의 메타이론적 기반을 비판적 실재론의 층화된 존재론에 둔다. \citet{Bhaskar1975}에 따르면 세계는 세 층위로 구성된다. 실재적 영역은 관찰 여부와 무관하게 존재하는 구조, 메커니즘, 인과적 권능을 포함한다. 실제적 영역은 이러한 구조와 메커니즘이 특정 조건에서 작동하여 사건을 발생시키는 영역이다. 경험적 영역은 발생한 사건 가운데 관찰 가능한 일부를 가리킨다. 실제적 영역의 모든 사건이 경험되는 것은 아니며, 관찰된 현상이 실재적 영역의 전부를 나타내는 것도 아니다.

이 존재론적 구분은 역행추론(retroduction)의 논리를 수반한다~\citep{Bhaskar1975}. 역행추론은 관찰된 경험적 사건으로부터, 그 사건을 가능하게 한 구조와 메커니즘을 역으로 추론하는 방법이다. 귀납이나 연역과 구분되며, \enquote{이 현상이 가능하려면 어떤 구조가 존재해야 하는가}를 묻는다.

\subsection{Seoul}
\label{subsec:2-2}

\subsubsection{도시로서의 서울}
\label{subsubsec:2-2-1}

대한민국의 수도인 서울은 국가 전체 국토 면적의 0.6\%에 인구의 20\%가 거주하는 거대한 밀집 도시이다. 높은 밀집도와 기술력의 집중에 따라, 서울은 대한민국 디지털 전환의 최전선에 존재해왔다. 그러나 서울의 내부를 면밀히 들여다보면 각 25개 자치구의 인구 구조와 경제적 자생력에 따른 유의미한 격차가 존재한다. 본 절에서는 서울 내 각 자치구별 특징을 조사하고, 서울의 인구 지형의 특징에 대해서 논의하고자 한다.

최근 서울의 인구 지형은 급격한 고령화와 1인 가구의 팽창이라는 두 가지 축으로 재편되고 있다. 강북구(26.03\%)와 도봉구(25.85\%) 등 서울 북부 외곽 지역은 고령 인구 비율이 이미 25\%를 상회하며 초고령사회적 특성을 띠는 반면, 강남권과 주요 대학가 주변은 상대적으로 젊은 인구 구조를 유지하고 있어 지역 간 연령 분포의 불균형이 뚜렷하다. 이러한 인구 통계적 차이는 각 구의 평균 소득 수준 및 경제적 활력과 결합하여, 본 연구 \ref{subsec:3-1}절에서 확인된 가격 부담(Affordability)이나 디지털 기기 보유(Devices) 역량의 차이를 발생시키는 구조적 배경이 된다.

\fig{image4.png}{서울 25개 자치구 인구\textperiodcentered{}연령 구조}{fig:2-pop}

또한, 서울의 정책 시행 체계는 서울특별시라는 광역 자치단체와 25개 기초 자치단체 사이의 유기적인 협력을 기반으로 작동한다. 정책의 큰 줄기는 서울시 본청에서 수립하는 하향식(Top-down) 방식을 따르는데, 이는 디지털 포용을 위한 광역 단위의 인프라 구축이나 표준 가이드라인 설정에 핵심적인 역할을 한다. 동시에 서울시는 시민참여예산제와 주민자치회 등을 통해 지역 주민이 실생활에서 느끼는 디지털 불편 사항을 정책에 직접 반영하는 상향식(Bottom-up) 의사결정 구조도 함께 운용하고 있다. 자치구는 이러한 광역의 계획을 바탕으로 지역 내 의회 심의와 조례 제정을 거쳐 예산을 확정하며, 최종적으로 동 주민센터를 통해 시민들에게 밀착형 서비스를 전달한다. 결국 서울의 디지털 격차 해소 정책은 중앙의 일괄적인 집행이 아닌, 각 자치구의 행정 역량과 시민 참여가 결합하여 구체화되는 다층적인 과정을 거쳐 실현된다고 볼 수 있다.

\subsubsection{연결성 격차 현황과 관련 법령 및 정책 현황}
\label{subsubsec:2-2-2}

서울의 디지털 연결성 실태는 자치구별 무선국 인프라 밀도의 편차를 통해 민간 주도 인프라 구축의 한계를 여실히 드러낸다. 전파누리의 2020--2025년 무선국 통계에 따르면, 서울의 통신 품질을 결정짓는 핵심 인프라인 '기지국'과 '이동중계국'은 수익성이 높고 트래픽이 집중되는 특정 지역에 압도적으로 쏠려 있다. 강남구(기지국 19,190개, 이동중계국 2,711개)와 송파구(기지국 14,751개) 등 동남권 지역은 촘촘한 망을 보유한 반면, 도봉구(기지국 4,008개)나 강북구(기지국 4,700개) 등 북부 외곽 지역은 기지국 수에서 최대 4.8배의 격차를 보인다. 특히 소규모 사업장과 사물인터넷(IoT) 연결에 활용되는 '간이무선국'의 경우 강남구(24,236개)가 도봉구(1,224개)보다 약 20배나 많아, 자치구 간 디지털 생태계의 부익부 빈익빈 현상이 물리적 인프라 수준에서 이미 고착화되고 있음을 시사한다.

\fig{image5.png}{자치구별 무선국 인프라 격차 (전파누리 2020--2025)}{fig:2-radio}

이러한 민간 인프라의 지역적 불균형을 교정하고 시민의 통신기본권을 보장하기 위해, 서울시는 행정적 대안으로서 직접 '기간통신사업자'로 등판하는 정책적 결단을 내렸다. 2025년 전기통신사업법 개정에 따라 지자체 최초로 통신사업자 지위를 확보한 서울시는, 민간 통신사가 수익성 문제로 투자를 미진하게 하는 소외 지역에 5,000km 이상의 '자가 광케이블망'을 직접 연결하여 공공 와이파이와 IoT 인프라를 확충하고 있다. 현재 서울시는 전국 최대 규모인 3만4천여 대의 공공와이파이를 운영하고 있으며, 이는 거리, 공원, 시장, 복지시설, 시내버스 등에 폭넓게 설치돼 연간 약 2천억 원의 시민 통신비 절감 효과를 거두고 있다. 특히 시내버스 7천여 대에 설치된 이동형 와이파이는 시민들이 일상에서 가장 체감하는 디지털 인프라로 자리잡고 있다. 서울시는 공공와이파이 확산을 통해 정보 접근의 격차를 줄이고, 매년 수천억 원에 달하는 가계 통신비 부담을 덜겠다는 계획이며, 또한 민간 임대망 의존도를 낮춰 확보된 재원을 시민 편익 사업에 재투자함으로써, 지속 가능한 선순환 구조를 마련하고자 하고 있다.

이러한 인프라 격차를 해소하기 위한 정책적 대응은 물리적 망 확충을 넘어, 시민의 실질적인 비용 부담을 낮추는 '경제적 접근성(Affordability)' 강화에 집중되고 있다. 정부와 지자체는 이용자가 자신의 통신 소비 패턴을 파악하고 합리적인 선택을 할 수 있도록 돕는 '스마트초이스(Smart Choice)' 플랫폼을 운영하고 있다. 스마트초이스는 단말기 지원금 조회, 맞춤형 요금제 추천, 미환급액 조회 등 통합 서비스를 제공하여 정보 비대칭성을 해소하고 가계 통신비 절감을 유도한다.

특히 최근 단행된 5G 요금제 개편은 기존 고가 요금제 중심의 시장 구조를 중저가 중심으로 전환했다는 점에서 큰 의의를 갖는다. 과거 4\,$\sim$\,5만 원대에 머물던 5G 최저 요금 구간을 3만 원대로 낮추었으며, 실제 데이터 사용량에 부합하도록 30\,$\sim$\,100GB 사이의 중간 요금제 구간을 촘촘하게 신설하였다. 이러한 개편은 이용자가 쓰지도 않는 데이터에 높은 비용을 지불하던 구조적 모순을 해결하고, 2024년 기준 600만 명 이상의 가입자가 중저가 요금제로 전환하는 실질적인 정책 성과로 이어졌다.

취약계층 및 특정 연령대를 위한 맞춤형 전용 요금제 역시 대폭 강화되었다. 대한민국 대표 이동통신사 3사에서는 청년층(만 34세 이하)을 위해, 데이터 수요가 많은 특성을 고려하여 일반 요금제와 동일한 가격에 데이터 제공량을 최대 2배까지 확대하고, 가입 연령 상한을 기존 만 29세에서 34세로 상향하여 수혜 대상을 넓혔다. 또한 LG U+에서는 고령층(만 65세 이상)을 위해, 상대적으로 데이터 사용량이 적고 가격 민감도가 높은 점을 반영하여, 일반 요금제 대비 최대 20\% 저렴한 전용 요금제를 연령대별(65세\textperiodcentered{}70세\textperiodcentered{}80세 등)로 세분화하여 출시하였다. 나아가 온라인 전용 요금제를 통해 비대면 가입 활성화 추세에 맞춰 대리점 방문 없이 온라인 채널로 가입 시 일반 요금제 대비 약 30\% 저렴하면서 약정이 없는 무약정 요금제를 도입하여 선택권을 넓혔다.

결과적으로 서울시는 스마트초이스를 통한 정보의 투명성, 통신 3사의 합리적 요금 체계, 그리고 서울시의 직접적인 공공 인프라라는 삼각 편대를 통해 모든 시민의 보편적 디지털 기본권을 보장하는 다각적인 포용 정책을 완성해 나가고 있다.

이에 더불어, 2026년 서울시 행정 서비스의 가장 큰 변화는 모바일 플랫폼의 단일화와 고도화이다. 서울시는 기존에 분절되어 운영되던 '서울시민카드'와 '서울지갑' 서비스를 2026년 3월 31일부로 종료하고, 통합 모바일 플랫폼인 '서울온(Seoul ON)'으로 기능을 일원화한다. 이는 시민들이 도서관, 체육시설 등 공공시설 이용부터 장애인\textperiodcentered{}국가유공자 자격 확인까지 하나의 앱에서 해결할 수 있게 함으로써, 서비스 파편화로 인한 디지털 행정의 진입장벽을 낮추고한 것이다. 특히 '서울온'은 공공 마이데이터와 연계하여 별도의 서류 제출 없이도 취약계층 자격을 실시간 확인하는 기능을 강화하여 행정 효율성을 높이는 기능을 한다.

나아가 서울시는 인프라 확충과 서비스 통합이 시민의 실질적인 이용 역량 강화로 이어질 수 있도록 'AI디지털배움터' 중심의 통합 포용 정책을 추진하고 있다. 서울 내 10개 거점을 포함하여, 전국 단위의 디지털 배움터 거점을 활용해 고령층 및 디지털 취약계층을 대상으로 한 대면 교육을 상시 운영하고 있으며, 단순한 인터넷 활용 교육을 넘어 실생활 중심의 AI 체험 교육과 디지털 역기능 예방 교육을 통합하여 제공한다. 특히 거동이 불편한 중증 장애인 등을 위해 '찾아가는 1:1 방문 교육'을 확대하고, 온\textperiodcentered{}오프라인이 연계된 수준별 맞춤형 교육 체계를 구축함으로써 교육 사각지대를 해소하고 있다. 이는 하드웨어(인프라)와 소프트웨어(행정 서비스)가 갖춰지더라도 이를 활용할 인적 역량이 부족하면 격차가 해소되지 않는다는 판단에 따른 것이다. 이처럼 서울시는 법령에 기반한 인프라 안정성 확보, '서울온'을 통한 서비스 전달 체계의 혁신, 그리고 교육을 통한 이용자 역량 제고라는 삼각 편대를 통해 보편적이고 의미 있는 연결성(UMC)을 실질적으로 구현해 나가고 있다.
)
\section{이론적 기반과 해석적 맥락}
\label{sec:theory}

\noindent{\sffamily\small\color{muted}\textbf{CHAPTER 2}}

\begin{quoteblock}
사람 기반과 장소 기반 접근을 통합하는 장소민감적(place-sensitive) 접근의 학술\textperiodcentered{}정책적 근거를 정리하고, 서울의 인구\textperiodcentered{}인프라\textperiodcentered{}정책 환경을 분석의 해석적 맥락으로 제시한다.
\end{quoteblock}

본 장은 분석에 관한 이론적 기반과 3장의 해석적 맥락을 제시한다. 먼저, 디지털 격차에 영향을 미치는 개인적 요인과 지역적 요인을 정리한다. 이후, 공공정책에서 사람 기반 접근(people-based)과 장소 기반 접근(place-based)의 구분과 최근 장소민감적 접근의 논의배경을 검토한다. 그 다음으로, 본 프로젝트의 메타이론적 틀인 비판적 실재론과 실제 연구 적용 사례를 검토한다. 마지막으로 서울의 인구, 행정, 인프라 조건을 검토해 분석 결과의 해석에 필요한 맥락 정보를 제공한다.

\subsection{Policy}
\label{subsec:2-1}

공공정책의 설계에서 개입의 대상을 사람의 특성에 두는 접근(people-based approach)과 장소의 조건에 두는 접근(place-based approach)의 구분은 오래된 논쟁이다. 사람 기반 정책은 개인의 역량, 소득, 교육 등 인적 특성을 직접 지원함으로써 결과를 개선하고자 한다. 장소 기반 정책은 특정 지역의 인프라, 제도, 환경 조건을 개선함으로써 그 지역 거주자의 결과를 간접적으로 향상시키고자 한다~\citep{NeumarkSimpson2015}. 두 접근은 상호 배타적이지 않으나, 정책 자원의 배분과 개입의 경로가 상이하므로 정책 수행의 관점으로 분리해 설계될 필요가 있다.

\subsubsection{디지털 격차의 영향 요인: 개인 특성과 지역 조건}
\label{subsubsec:2-1-1}

디지털 격차 문헌에서 개인 특성에 주목하는 연구는, 동일 지역 내에서도 개인 특성이 디지털 활용 격차를 설명하는 주된 요인임을 보여준다. \citet{Hargittai2010}는 인터넷 접근 이후의 격차, 즉 활용 방식의 차이가 실질적 불평등을 결정한다고 제시하였다. 사회경제적 지위에 따라 디지털 활용의 격차가 지속될 수 있음을 시사하는 결과이다. \citet{ChipevaEtAl2018}은 다양한 국가에서 실용성이 기술 사용의 가장 강력한 동기임을 보이는 한편, 개인의 내적 특성이 디지털 기기 사용에 유의미한 영향을 미친다는 점을 제시하였다. \citet{BorgesEtAl2020}은 앙골라 사례를 통해 기기 보급 이전에 습관과 ICT 기술 숙련도가 격차의 핵심 요인임을 보고하였다. \citet{Hilbert2011}의 라틴아메리카 18개국 분석은 남녀 간 디지털 격차가 존재하며, 국가 수준의 젠더 평등 지수가 높을수록 인터넷 사용이 활발함을 보였다.

한편, 개인 특성을 통제한 이후에도 지역 환경이 디지털 활용에 대해 추가적 설명력을 갖는다는 증거도 축적되어 있다. \citet{MossbergerEtAl2006}은 미국을 대상으로 HLM을 적용하여, 인터넷 접근 격차가 인종 자체보다 고빈곤 지역 환경에 기인함을 보였다. 이는 지역의 구조적 조건이 개인의 결과에 독립적 영향을 미칠 수 있음을 시사한다. \citet{CrangEtAl2006}은 영국 뉴캐슬의 질적 사례 비교를 통해, 지역 환경이 디지털 네트워크의 사용 방식 자체를 구조화함을 제시하였다. \citet{HongEtAl2017}의 중국 45세 이상 성인 분석은, 인터넷 이용자가 고소득, 고학력, 도시 거주 집단에 집중되어 있음을 보여, 지역 수준의 인프라와 기회 구조가 이용 가능성을 선별할 수 있음을 시사하였다.

\subsubsection{정책적 구분: 사람 기반 접근, 장소 기반 접근, 장소민감적 접근}
\label{subsubsec:2-1-2}

앞 절에서 검토한 두 유형의 요인은 공공정책 설계에서 사람 기반 접근(people-based approach)과 장소 기반 접근(place-based approach)의 구분에 대응한다. 사람 기반 정책은 개인의 역량, 소득, 교육 등 인적 특성을 직접 지원함으로써 결과를 개선하고자 한다. 장소 기반 정책은 특정 지역의 인프라, 제도, 환경 조건을 개선함으로써 그 지역 거주자의 결과를 간접적으로 향상시키고자 한다~\citep{NeumarkSimpson2015}. 두 접근은 상호 배타적이지 않으나, 정책 자원의 배분과 개입의 경로가 상이하므로 정책 수행의 관점에서 분리하여 설계될 필요가 있다.

그러나 두 접근 사이의 관계는 단순한 병렬이 아니다. \citet{Hindman2000}의 미국 시계열 분석은 소득, 연령, 교육이 지리보다 강력한 예측변수임을 보이면서도, 시간 경과에 따라 격차가 고착화되는 패턴을 보고하였다. 이는 개인 특성과 지역 구조가 시간 속에서 상호작용할 수 있음을 의미한다. \citet{vanDijk2020}는 동기 결여에서 물질적 접근, 기술 능력, 실질적 활용으로 이어지는 순차적 모형을 제시하여, 개인 요인과 구조 요인이 누적적으로 작용함을 제안하였다.

이러한 맥락에서, 최근 정책 및 학술 문헌은 사람 기반 또는 장소 기반의 양자택일이 아닌, 장소의 맥락 안에서 사람의 경험을 이해하는 장소민감적(place-sensitive) 접근으로 이행하고 있다. 생활 경험(lived experience)과 장소 기반 통찰(place-based insight)을 정량적 증거와 동등한 지식으로 인정하려는 시도가 그 예이다~\citep{OECD2025}. 이 접근은 사람 기반 요인과 장소 기반 요인을 분석적으로 분리하되, 두 요인이 특정 장소에서 어떻게 결합하는지를 함께 고려할 것을 요구한다.

\subsubsection{메타이론적 틀: 비판적 실재론의 층화된 존재론}
\label{subsubsec:2-1-3}

본 프로젝트는 사람 기반 요인과 장소 기반 요인의 분리를 분석 설계의 기본 원칙으로 채택하며, 그 설계의 메타이론적 기반을 비판적 실재론의 층화된 존재론에 둔다. \citet{Bhaskar1975}에 따르면 세계는 세 층위로 구성된다. 실재적 영역은 관찰 여부와 무관하게 존재하는 구조, 메커니즘, 인과적 권능을 포함한다. 실제적 영역은 이러한 구조와 메커니즘이 특정 조건에서 작동하여 사건을 발생시키는 영역이다. 경험적 영역은 발생한 사건 가운데 관찰 가능한 일부를 가리킨다. 실제적 영역의 모든 사건이 경험되는 것은 아니며, 관찰된 현상이 실재적 영역의 전부를 나타내는 것도 아니다.

이 존재론적 구분은 역행추론(retroduction)의 논리를 수반한다~\citep{Bhaskar1975}. 역행추론은 관찰된 경험적 사건으로부터, 그 사건을 가능하게 한 구조와 메커니즘을 역으로 추론하는 방법이다. 귀납이나 연역과 구분되며, \enquote{이 현상이 가능하려면 어떤 구조가 존재해야 하는가}를 묻는다.

\subsection{Seoul}
\label{subsec:2-2}

\subsubsection{도시로서의 서울}
\label{subsubsec:2-2-1}

대한민국의 수도인 서울은 국가 전체 국토 면적의 0.6\%에 인구의 20\%가 거주하는 거대한 밀집 도시이다. 높은 밀집도와 기술력의 집중에 따라, 서울은 대한민국 디지털 전환의 최전선에 존재해왔다. 그러나 서울의 내부를 면밀히 들여다보면 각 25개 자치구의 인구 구조와 경제적 자생력에 따른 유의미한 격차가 존재한다. 본 절에서는 서울 내 각 자치구별 특징을 조사하고, 서울의 인구 지형의 특징에 대해서 논의하고자 한다.

최근 서울의 인구 지형은 급격한 고령화와 1인 가구의 팽창이라는 두 가지 축으로 재편되고 있다. 강북구(26.03\%)와 도봉구(25.85\%) 등 서울 북부 외곽 지역은 고령 인구 비율이 이미 25\%를 상회하며 초고령사회적 특성을 띠는 반면, 강남권과 주요 대학가 주변은 상대적으로 젊은 인구 구조를 유지하고 있어 지역 간 연령 분포의 불균형이 뚜렷하다. 이러한 인구 통계적 차이는 각 구의 평균 소득 수준 및 경제적 활력과 결합하여, 본 연구 \ref{subsec:3-1}절에서 확인된 가격 부담(Affordability)이나 디지털 기기 보유(Devices) 역량의 차이를 발생시키는 구조적 배경이 된다.

\fig{image4.png}{서울 25개 자치구 인구\textperiodcentered{}연령 구조}{fig:2-pop}

또한, 서울의 정책 시행 체계는 서울특별시라는 광역 자치단체와 25개 기초 자치단체 사이의 유기적인 협력을 기반으로 작동한다. 정책의 큰 줄기는 서울시 본청에서 수립하는 하향식(Top-down) 방식을 따르는데, 이는 디지털 포용을 위한 광역 단위의 인프라 구축이나 표준 가이드라인 설정에 핵심적인 역할을 한다. 동시에 서울시는 시민참여예산제와 주민자치회 등을 통해 지역 주민이 실생활에서 느끼는 디지털 불편 사항을 정책에 직접 반영하는 상향식(Bottom-up) 의사결정 구조도 함께 운용하고 있다. 자치구는 이러한 광역의 계획을 바탕으로 지역 내 의회 심의와 조례 제정을 거쳐 예산을 확정하며, 최종적으로 동 주민센터를 통해 시민들에게 밀착형 서비스를 전달한다. 결국 서울의 디지털 격차 해소 정책은 중앙의 일괄적인 집행이 아닌, 각 자치구의 행정 역량과 시민 참여가 결합하여 구체화되는 다층적인 과정을 거쳐 실현된다고 볼 수 있다.

\subsubsection{연결성 격차 현황과 관련 법령 및 정책 현황}
\label{subsubsec:2-2-2}

서울의 디지털 연결성 실태는 자치구별 무선국 인프라 밀도의 편차를 통해 민간 주도 인프라 구축의 한계를 여실히 드러낸다. 전파누리의 2020--2025년 무선국 통계에 따르면, 서울의 통신 품질을 결정짓는 핵심 인프라인 '기지국'과 '이동중계국'은 수익성이 높고 트래픽이 집중되는 특정 지역에 압도적으로 쏠려 있다. 강남구(기지국 19,190개, 이동중계국 2,711개)와 송파구(기지국 14,751개) 등 동남권 지역은 촘촘한 망을 보유한 반면, 도봉구(기지국 4,008개)나 강북구(기지국 4,700개) 등 북부 외곽 지역은 기지국 수에서 최대 4.8배의 격차를 보인다. 특히 소규모 사업장과 사물인터넷(IoT) 연결에 활용되는 '간이무선국'의 경우 강남구(24,236개)가 도봉구(1,224개)보다 약 20배나 많아, 자치구 간 디지털 생태계의 부익부 빈익빈 현상이 물리적 인프라 수준에서 이미 고착화되고 있음을 시사한다.

\fig{image5.png}{자치구별 무선국 인프라 격차 (전파누리 2020--2025)}{fig:2-radio}

이러한 민간 인프라의 지역적 불균형을 교정하고 시민의 통신기본권을 보장하기 위해, 서울시는 행정적 대안으로서 직접 '기간통신사업자'로 등판하는 정책적 결단을 내렸다. 2025년 전기통신사업법 개정에 따라 지자체 최초로 통신사업자 지위를 확보한 서울시는, 민간 통신사가 수익성 문제로 투자를 미진하게 하는 소외 지역에 5,000km 이상의 '자가 광케이블망'을 직접 연결하여 공공 와이파이와 IoT 인프라를 확충하고 있다. 현재 서울시는 전국 최대 규모인 3만4천여 대의 공공와이파이를 운영하고 있으며, 이는 거리, 공원, 시장, 복지시설, 시내버스 등에 폭넓게 설치돼 연간 약 2천억 원의 시민 통신비 절감 효과를 거두고 있다. 특히 시내버스 7천여 대에 설치된 이동형 와이파이는 시민들이 일상에서 가장 체감하는 디지털 인프라로 자리잡고 있다. 서울시는 공공와이파이 확산을 통해 정보 접근의 격차를 줄이고, 매년 수천억 원에 달하는 가계 통신비 부담을 덜겠다는 계획이며, 또한 민간 임대망 의존도를 낮춰 확보된 재원을 시민 편익 사업에 재투자함으로써, 지속 가능한 선순환 구조를 마련하고자 하고 있다.

이러한 인프라 격차를 해소하기 위한 정책적 대응은 물리적 망 확충을 넘어, 시민의 실질적인 비용 부담을 낮추는 '경제적 접근성(Affordability)' 강화에 집중되고 있다. 정부와 지자체는 이용자가 자신의 통신 소비 패턴을 파악하고 합리적인 선택을 할 수 있도록 돕는 '스마트초이스(Smart Choice)' 플랫폼을 운영하고 있다. 스마트초이스는 단말기 지원금 조회, 맞춤형 요금제 추천, 미환급액 조회 등 통합 서비스를 제공하여 정보 비대칭성을 해소하고 가계 통신비 절감을 유도한다.

특히 최근 단행된 5G 요금제 개편은 기존 고가 요금제 중심의 시장 구조를 중저가 중심으로 전환했다는 점에서 큰 의의를 갖는다. 과거 4\,$\sim$\,5만 원대에 머물던 5G 최저 요금 구간을 3만 원대로 낮추었으며, 실제 데이터 사용량에 부합하도록 30\,$\sim$\,100GB 사이의 중간 요금제 구간을 촘촘하게 신설하였다. 이러한 개편은 이용자가 쓰지도 않는 데이터에 높은 비용을 지불하던 구조적 모순을 해결하고, 2024년 기준 600만 명 이상의 가입자가 중저가 요금제로 전환하는 실질적인 정책 성과로 이어졌다.

취약계층 및 특정 연령대를 위한 맞춤형 전용 요금제 역시 대폭 강화되었다. 대한민국 대표 이동통신사 3사에서는 청년층(만 34세 이하)을 위해, 데이터 수요가 많은 특성을 고려하여 일반 요금제와 동일한 가격에 데이터 제공량을 최대 2배까지 확대하고, 가입 연령 상한을 기존 만 29세에서 34세로 상향하여 수혜 대상을 넓혔다. 또한 LG U+에서는 고령층(만 65세 이상)을 위해, 상대적으로 데이터 사용량이 적고 가격 민감도가 높은 점을 반영하여, 일반 요금제 대비 최대 20\% 저렴한 전용 요금제를 연령대별(65세\textperiodcentered{}70세\textperiodcentered{}80세 등)로 세분화하여 출시하였다. 나아가 온라인 전용 요금제를 통해 비대면 가입 활성화 추세에 맞춰 대리점 방문 없이 온라인 채널로 가입 시 일반 요금제 대비 약 30\% 저렴하면서 약정이 없는 무약정 요금제를 도입하여 선택권을 넓혔다.

결과적으로 서울시는 스마트초이스를 통한 정보의 투명성, 통신 3사의 합리적 요금 체계, 그리고 서울시의 직접적인 공공 인프라라는 삼각 편대를 통해 모든 시민의 보편적 디지털 기본권을 보장하는 다각적인 포용 정책을 완성해 나가고 있다.

이에 더불어, 2026년 서울시 행정 서비스의 가장 큰 변화는 모바일 플랫폼의 단일화와 고도화이다. 서울시는 기존에 분절되어 운영되던 '서울시민카드'와 '서울지갑' 서비스를 2026년 3월 31일부로 종료하고, 통합 모바일 플랫폼인 '서울온(Seoul ON)'으로 기능을 일원화한다. 이는 시민들이 도서관, 체육시설 등 공공시설 이용부터 장애인\textperiodcentered{}국가유공자 자격 확인까지 하나의 앱에서 해결할 수 있게 함으로써, 서비스 파편화로 인한 디지털 행정의 진입장벽을 낮추고한 것이다. 특히 '서울온'은 공공 마이데이터와 연계하여 별도의 서류 제출 없이도 취약계층 자격을 실시간 확인하는 기능을 강화하여 행정 효율성을 높이는 기능을 한다.

나아가 서울시는 인프라 확충과 서비스 통합이 시민의 실질적인 이용 역량 강화로 이어질 수 있도록 'AI디지털배움터' 중심의 통합 포용 정책을 추진하고 있다. 서울 내 10개 거점을 포함하여, 전국 단위의 디지털 배움터 거점을 활용해 고령층 및 디지털 취약계층을 대상으로 한 대면 교육을 상시 운영하고 있으며, 단순한 인터넷 활용 교육을 넘어 실생활 중심의 AI 체험 교육과 디지털 역기능 예방 교육을 통합하여 제공한다. 특히 거동이 불편한 중증 장애인 등을 위해 '찾아가는 1:1 방문 교육'을 확대하고, 온\textperiodcentered{}오프라인이 연계된 수준별 맞춤형 교육 체계를 구축함으로써 교육 사각지대를 해소하고 있다. 이는 하드웨어(인프라)와 소프트웨어(행정 서비스)가 갖춰지더라도 이를 활용할 인적 역량이 부족하면 격차가 해소되지 않는다는 판단에 따른 것이다. 이처럼 서울시는 법령에 기반한 인프라 안정성 확보, '서울온'을 통한 서비스 전달 체계의 혁신, 그리고 교육을 통한 이용자 역량 제고라는 삼각 편대를 통해 보편적이고 의미 있는 연결성(UMC)을 실질적으로 구현해 나가고 있다.
)
\section{이론적 기반과 해석적 맥락}
\label{sec:theory}

\noindent{\sffamily\small\color{muted}\textbf{CHAPTER 2}}

\begin{quoteblock}
사람 기반과 장소 기반 접근을 통합하는 장소민감적(place-sensitive) 접근의 학술\textperiodcentered{}정책적 근거를 정리하고, 서울의 인구\textperiodcentered{}인프라\textperiodcentered{}정책 환경을 분석의 해석적 맥락으로 제시한다.
\end{quoteblock}

본 장은 분석에 관한 이론적 기반과 3장의 해석적 맥락을 제시한다. 먼저, 디지털 격차에 영향을 미치는 개인적 요인과 지역적 요인을 정리한다. 이후, 공공정책에서 사람 기반 접근(people-based)과 장소 기반 접근(place-based)의 구분과 최근 장소민감적 접근의 논의배경을 검토한다. 그 다음으로, 본 프로젝트의 메타이론적 틀인 비판적 실재론과 실제 연구 적용 사례를 검토한다. 마지막으로 서울의 인구, 행정, 인프라 조건을 검토해 분석 결과의 해석에 필요한 맥락 정보를 제공한다.

\subsection{Policy}
\label{subsec:2-1}

공공정책의 설계에서 개입의 대상을 사람의 특성에 두는 접근(people-based approach)과 장소의 조건에 두는 접근(place-based approach)의 구분은 오래된 논쟁이다. 사람 기반 정책은 개인의 역량, 소득, 교육 등 인적 특성을 직접 지원함으로써 결과를 개선하고자 한다. 장소 기반 정책은 특정 지역의 인프라, 제도, 환경 조건을 개선함으로써 그 지역 거주자의 결과를 간접적으로 향상시키고자 한다~\citep{NeumarkSimpson2015}. 두 접근은 상호 배타적이지 않으나, 정책 자원의 배분과 개입의 경로가 상이하므로 정책 수행의 관점으로 분리해 설계될 필요가 있다.

\subsubsection{디지털 격차의 영향 요인: 개인 특성과 지역 조건}
\label{subsubsec:2-1-1}

디지털 격차 문헌에서 개인 특성에 주목하는 연구는, 동일 지역 내에서도 개인 특성이 디지털 활용 격차를 설명하는 주된 요인임을 보여준다. \citet{Hargittai2010}는 인터넷 접근 이후의 격차, 즉 활용 방식의 차이가 실질적 불평등을 결정한다고 제시하였다. 사회경제적 지위에 따라 디지털 활용의 격차가 지속될 수 있음을 시사하는 결과이다. \citet{ChipevaEtAl2018}은 다양한 국가에서 실용성이 기술 사용의 가장 강력한 동기임을 보이는 한편, 개인의 내적 특성이 디지털 기기 사용에 유의미한 영향을 미친다는 점을 제시하였다. \citet{BorgesEtAl2020}은 앙골라 사례를 통해 기기 보급 이전에 습관과 ICT 기술 숙련도가 격차의 핵심 요인임을 보고하였다. \citet{Hilbert2011}의 라틴아메리카 18개국 분석은 남녀 간 디지털 격차가 존재하며, 국가 수준의 젠더 평등 지수가 높을수록 인터넷 사용이 활발함을 보였다.

한편, 개인 특성을 통제한 이후에도 지역 환경이 디지털 활용에 대해 추가적 설명력을 갖는다는 증거도 축적되어 있다. \citet{MossbergerEtAl2006}은 미국을 대상으로 HLM을 적용하여, 인터넷 접근 격차가 인종 자체보다 고빈곤 지역 환경에 기인함을 보였다. 이는 지역의 구조적 조건이 개인의 결과에 독립적 영향을 미칠 수 있음을 시사한다. \citet{CrangEtAl2006}은 영국 뉴캐슬의 질적 사례 비교를 통해, 지역 환경이 디지털 네트워크의 사용 방식 자체를 구조화함을 제시하였다. \citet{HongEtAl2017}의 중국 45세 이상 성인 분석은, 인터넷 이용자가 고소득, 고학력, 도시 거주 집단에 집중되어 있음을 보여, 지역 수준의 인프라와 기회 구조가 이용 가능성을 선별할 수 있음을 시사하였다.

\subsubsection{정책적 구분: 사람 기반 접근, 장소 기반 접근, 장소민감적 접근}
\label{subsubsec:2-1-2}

앞 절에서 검토한 두 유형의 요인은 공공정책 설계에서 사람 기반 접근(people-based approach)과 장소 기반 접근(place-based approach)의 구분에 대응한다. 사람 기반 정책은 개인의 역량, 소득, 교육 등 인적 특성을 직접 지원함으로써 결과를 개선하고자 한다. 장소 기반 정책은 특정 지역의 인프라, 제도, 환경 조건을 개선함으로써 그 지역 거주자의 결과를 간접적으로 향상시키고자 한다~\citep{NeumarkSimpson2015}. 두 접근은 상호 배타적이지 않으나, 정책 자원의 배분과 개입의 경로가 상이하므로 정책 수행의 관점에서 분리하여 설계될 필요가 있다.

그러나 두 접근 사이의 관계는 단순한 병렬이 아니다. \citet{Hindman2000}의 미국 시계열 분석은 소득, 연령, 교육이 지리보다 강력한 예측변수임을 보이면서도, 시간 경과에 따라 격차가 고착화되는 패턴을 보고하였다. 이는 개인 특성과 지역 구조가 시간 속에서 상호작용할 수 있음을 의미한다. \citet{vanDijk2020}는 동기 결여에서 물질적 접근, 기술 능력, 실질적 활용으로 이어지는 순차적 모형을 제시하여, 개인 요인과 구조 요인이 누적적으로 작용함을 제안하였다.

이러한 맥락에서, 최근 정책 및 학술 문헌은 사람 기반 또는 장소 기반의 양자택일이 아닌, 장소의 맥락 안에서 사람의 경험을 이해하는 장소민감적(place-sensitive) 접근으로 이행하고 있다. 생활 경험(lived experience)과 장소 기반 통찰(place-based insight)을 정량적 증거와 동등한 지식으로 인정하려는 시도가 그 예이다~\citep{OECD2025}. 이 접근은 사람 기반 요인과 장소 기반 요인을 분석적으로 분리하되, 두 요인이 특정 장소에서 어떻게 결합하는지를 함께 고려할 것을 요구한다.

\subsubsection{메타이론적 틀: 비판적 실재론의 층화된 존재론}
\label{subsubsec:2-1-3}

본 프로젝트는 사람 기반 요인과 장소 기반 요인의 분리를 분석 설계의 기본 원칙으로 채택하며, 그 설계의 메타이론적 기반을 비판적 실재론의 층화된 존재론에 둔다. \citet{Bhaskar1975}에 따르면 세계는 세 층위로 구성된다. 실재적 영역은 관찰 여부와 무관하게 존재하는 구조, 메커니즘, 인과적 권능을 포함한다. 실제적 영역은 이러한 구조와 메커니즘이 특정 조건에서 작동하여 사건을 발생시키는 영역이다. 경험적 영역은 발생한 사건 가운데 관찰 가능한 일부를 가리킨다. 실제적 영역의 모든 사건이 경험되는 것은 아니며, 관찰된 현상이 실재적 영역의 전부를 나타내는 것도 아니다.

이 존재론적 구분은 역행추론(retroduction)의 논리를 수반한다~\citep{Bhaskar1975}. 역행추론은 관찰된 경험적 사건으로부터, 그 사건을 가능하게 한 구조와 메커니즘을 역으로 추론하는 방법이다. 귀납이나 연역과 구분되며, \enquote{이 현상이 가능하려면 어떤 구조가 존재해야 하는가}를 묻는다.

\subsection{Seoul}
\label{subsec:2-2}

\subsubsection{도시로서의 서울}
\label{subsubsec:2-2-1}

대한민국의 수도인 서울은 국가 전체 국토 면적의 0.6\%에 인구의 20\%가 거주하는 거대한 밀집 도시이다. 높은 밀집도와 기술력의 집중에 따라, 서울은 대한민국 디지털 전환의 최전선에 존재해왔다. 그러나 서울의 내부를 면밀히 들여다보면 각 25개 자치구의 인구 구조와 경제적 자생력에 따른 유의미한 격차가 존재한다. 본 절에서는 서울 내 각 자치구별 특징을 조사하고, 서울의 인구 지형의 특징에 대해서 논의하고자 한다.

최근 서울의 인구 지형은 급격한 고령화와 1인 가구의 팽창이라는 두 가지 축으로 재편되고 있다. 강북구(26.03\%)와 도봉구(25.85\%) 등 서울 북부 외곽 지역은 고령 인구 비율이 이미 25\%를 상회하며 초고령사회적 특성을 띠는 반면, 강남권과 주요 대학가 주변은 상대적으로 젊은 인구 구조를 유지하고 있어 지역 간 연령 분포의 불균형이 뚜렷하다. 이러한 인구 통계적 차이는 각 구의 평균 소득 수준 및 경제적 활력과 결합하여, 본 연구 \ref{subsec:3-1}절에서 확인된 가격 부담(Affordability)이나 디지털 기기 보유(Devices) 역량의 차이를 발생시키는 구조적 배경이 된다.

\fig{image4.png}{서울 25개 자치구 인구\textperiodcentered{}연령 구조}{fig:2-pop}

또한, 서울의 정책 시행 체계는 서울특별시라는 광역 자치단체와 25개 기초 자치단체 사이의 유기적인 협력을 기반으로 작동한다. 정책의 큰 줄기는 서울시 본청에서 수립하는 하향식(Top-down) 방식을 따르는데, 이는 디지털 포용을 위한 광역 단위의 인프라 구축이나 표준 가이드라인 설정에 핵심적인 역할을 한다. 동시에 서울시는 시민참여예산제와 주민자치회 등을 통해 지역 주민이 실생활에서 느끼는 디지털 불편 사항을 정책에 직접 반영하는 상향식(Bottom-up) 의사결정 구조도 함께 운용하고 있다. 자치구는 이러한 광역의 계획을 바탕으로 지역 내 의회 심의와 조례 제정을 거쳐 예산을 확정하며, 최종적으로 동 주민센터를 통해 시민들에게 밀착형 서비스를 전달한다. 결국 서울의 디지털 격차 해소 정책은 중앙의 일괄적인 집행이 아닌, 각 자치구의 행정 역량과 시민 참여가 결합하여 구체화되는 다층적인 과정을 거쳐 실현된다고 볼 수 있다.

\subsubsection{연결성 격차 현황과 관련 법령 및 정책 현황}
\label{subsubsec:2-2-2}

서울의 디지털 연결성 실태는 자치구별 무선국 인프라 밀도의 편차를 통해 민간 주도 인프라 구축의 한계를 여실히 드러낸다. 전파누리의 2020--2025년 무선국 통계에 따르면, 서울의 통신 품질을 결정짓는 핵심 인프라인 '기지국'과 '이동중계국'은 수익성이 높고 트래픽이 집중되는 특정 지역에 압도적으로 쏠려 있다. 강남구(기지국 19,190개, 이동중계국 2,711개)와 송파구(기지국 14,751개) 등 동남권 지역은 촘촘한 망을 보유한 반면, 도봉구(기지국 4,008개)나 강북구(기지국 4,700개) 등 북부 외곽 지역은 기지국 수에서 최대 4.8배의 격차를 보인다. 특히 소규모 사업장과 사물인터넷(IoT) 연결에 활용되는 '간이무선국'의 경우 강남구(24,236개)가 도봉구(1,224개)보다 약 20배나 많아, 자치구 간 디지털 생태계의 부익부 빈익빈 현상이 물리적 인프라 수준에서 이미 고착화되고 있음을 시사한다.

\fig{image5.png}{자치구별 무선국 인프라 격차 (전파누리 2020--2025)}{fig:2-radio}

이러한 민간 인프라의 지역적 불균형을 교정하고 시민의 통신기본권을 보장하기 위해, 서울시는 행정적 대안으로서 직접 '기간통신사업자'로 등판하는 정책적 결단을 내렸다. 2025년 전기통신사업법 개정에 따라 지자체 최초로 통신사업자 지위를 확보한 서울시는, 민간 통신사가 수익성 문제로 투자를 미진하게 하는 소외 지역에 5,000km 이상의 '자가 광케이블망'을 직접 연결하여 공공 와이파이와 IoT 인프라를 확충하고 있다. 현재 서울시는 전국 최대 규모인 3만4천여 대의 공공와이파이를 운영하고 있으며, 이는 거리, 공원, 시장, 복지시설, 시내버스 등에 폭넓게 설치돼 연간 약 2천억 원의 시민 통신비 절감 효과를 거두고 있다. 특히 시내버스 7천여 대에 설치된 이동형 와이파이는 시민들이 일상에서 가장 체감하는 디지털 인프라로 자리잡고 있다. 서울시는 공공와이파이 확산을 통해 정보 접근의 격차를 줄이고, 매년 수천억 원에 달하는 가계 통신비 부담을 덜겠다는 계획이며, 또한 민간 임대망 의존도를 낮춰 확보된 재원을 시민 편익 사업에 재투자함으로써, 지속 가능한 선순환 구조를 마련하고자 하고 있다.

이러한 인프라 격차를 해소하기 위한 정책적 대응은 물리적 망 확충을 넘어, 시민의 실질적인 비용 부담을 낮추는 '경제적 접근성(Affordability)' 강화에 집중되고 있다. 정부와 지자체는 이용자가 자신의 통신 소비 패턴을 파악하고 합리적인 선택을 할 수 있도록 돕는 '스마트초이스(Smart Choice)' 플랫폼을 운영하고 있다. 스마트초이스는 단말기 지원금 조회, 맞춤형 요금제 추천, 미환급액 조회 등 통합 서비스를 제공하여 정보 비대칭성을 해소하고 가계 통신비 절감을 유도한다.

특히 최근 단행된 5G 요금제 개편은 기존 고가 요금제 중심의 시장 구조를 중저가 중심으로 전환했다는 점에서 큰 의의를 갖는다. 과거 4\,$\sim$\,5만 원대에 머물던 5G 최저 요금 구간을 3만 원대로 낮추었으며, 실제 데이터 사용량에 부합하도록 30\,$\sim$\,100GB 사이의 중간 요금제 구간을 촘촘하게 신설하였다. 이러한 개편은 이용자가 쓰지도 않는 데이터에 높은 비용을 지불하던 구조적 모순을 해결하고, 2024년 기준 600만 명 이상의 가입자가 중저가 요금제로 전환하는 실질적인 정책 성과로 이어졌다.

취약계층 및 특정 연령대를 위한 맞춤형 전용 요금제 역시 대폭 강화되었다. 대한민국 대표 이동통신사 3사에서는 청년층(만 34세 이하)을 위해, 데이터 수요가 많은 특성을 고려하여 일반 요금제와 동일한 가격에 데이터 제공량을 최대 2배까지 확대하고, 가입 연령 상한을 기존 만 29세에서 34세로 상향하여 수혜 대상을 넓혔다. 또한 LG U+에서는 고령층(만 65세 이상)을 위해, 상대적으로 데이터 사용량이 적고 가격 민감도가 높은 점을 반영하여, 일반 요금제 대비 최대 20\% 저렴한 전용 요금제를 연령대별(65세\textperiodcentered{}70세\textperiodcentered{}80세 등)로 세분화하여 출시하였다. 나아가 온라인 전용 요금제를 통해 비대면 가입 활성화 추세에 맞춰 대리점 방문 없이 온라인 채널로 가입 시 일반 요금제 대비 약 30\% 저렴하면서 약정이 없는 무약정 요금제를 도입하여 선택권을 넓혔다.

결과적으로 서울시는 스마트초이스를 통한 정보의 투명성, 통신 3사의 합리적 요금 체계, 그리고 서울시의 직접적인 공공 인프라라는 삼각 편대를 통해 모든 시민의 보편적 디지털 기본권을 보장하는 다각적인 포용 정책을 완성해 나가고 있다.

이에 더불어, 2026년 서울시 행정 서비스의 가장 큰 변화는 모바일 플랫폼의 단일화와 고도화이다. 서울시는 기존에 분절되어 운영되던 '서울시민카드'와 '서울지갑' 서비스를 2026년 3월 31일부로 종료하고, 통합 모바일 플랫폼인 '서울온(Seoul ON)'으로 기능을 일원화한다. 이는 시민들이 도서관, 체육시설 등 공공시설 이용부터 장애인\textperiodcentered{}국가유공자 자격 확인까지 하나의 앱에서 해결할 수 있게 함으로써, 서비스 파편화로 인한 디지털 행정의 진입장벽을 낮추고한 것이다. 특히 '서울온'은 공공 마이데이터와 연계하여 별도의 서류 제출 없이도 취약계층 자격을 실시간 확인하는 기능을 강화하여 행정 효율성을 높이는 기능을 한다.

나아가 서울시는 인프라 확충과 서비스 통합이 시민의 실질적인 이용 역량 강화로 이어질 수 있도록 'AI디지털배움터' 중심의 통합 포용 정책을 추진하고 있다. 서울 내 10개 거점을 포함하여, 전국 단위의 디지털 배움터 거점을 활용해 고령층 및 디지털 취약계층을 대상으로 한 대면 교육을 상시 운영하고 있으며, 단순한 인터넷 활용 교육을 넘어 실생활 중심의 AI 체험 교육과 디지털 역기능 예방 교육을 통합하여 제공한다. 특히 거동이 불편한 중증 장애인 등을 위해 '찾아가는 1:1 방문 교육'을 확대하고, 온\textperiodcentered{}오프라인이 연계된 수준별 맞춤형 교육 체계를 구축함으로써 교육 사각지대를 해소하고 있다. 이는 하드웨어(인프라)와 소프트웨어(행정 서비스)가 갖춰지더라도 이를 활용할 인적 역량이 부족하면 격차가 해소되지 않는다는 판단에 따른 것이다. 이처럼 서울시는 법령에 기반한 인프라 안정성 확보, '서울온'을 통한 서비스 전달 체계의 혁신, 그리고 교육을 통한 이용자 역량 제고라는 삼각 편대를 통해 보편적이고 의미 있는 연결성(UMC)을 실질적으로 구현해 나가고 있다.
)
\section{이론적 기반과 해석적 맥락}
\label{sec:theory}

\noindent{\sffamily\small\color{muted}\textbf{CHAPTER 2}}

\begin{quoteblock}
사람 기반과 장소 기반 접근을 통합하는 장소민감적(place-sensitive) 접근의 학술\textperiodcentered{}정책적 근거를 정리하고, 서울의 인구\textperiodcentered{}인프라\textperiodcentered{}정책 환경을 분석의 해석적 맥락으로 제시한다.
\end{quoteblock}

본 장은 분석에 관한 이론적 기반과 3장의 해석적 맥락을 제시한다. 먼저, 디지털 격차에 영향을 미치는 개인적 요인과 지역적 요인을 정리한다. 이후, 공공정책에서 사람 기반 접근(people-based)과 장소 기반 접근(place-based)의 구분과 최근 장소민감적 접근의 논의배경을 검토한다. 그 다음으로, 본 프로젝트의 메타이론적 틀인 비판적 실재론과 실제 연구 적용 사례를 검토한다. 마지막으로 서울의 인구, 행정, 인프라 조건을 검토해 분석 결과의 해석에 필요한 맥락 정보를 제공한다.

\subsection{Policy}
\label{subsec:2-1}

공공정책의 설계에서 개입의 대상을 사람의 특성에 두는 접근(people-based approach)과 장소의 조건에 두는 접근(place-based approach)의 구분은 오래된 논쟁이다. 사람 기반 정책은 개인의 역량, 소득, 교육 등 인적 특성을 직접 지원함으로써 결과를 개선하고자 한다. 장소 기반 정책은 특정 지역의 인프라, 제도, 환경 조건을 개선함으로써 그 지역 거주자의 결과를 간접적으로 향상시키고자 한다~\citep{NeumarkSimpson2015}. 두 접근은 상호 배타적이지 않으나, 정책 자원의 배분과 개입의 경로가 상이하므로 정책 수행의 관점으로 분리해 설계될 필요가 있다.

\subsubsection{디지털 격차의 영향 요인: 개인 특성과 지역 조건}
\label{subsubsec:2-1-1}

디지털 격차 문헌에서 개인 특성에 주목하는 연구는, 동일 지역 내에서도 개인 특성이 디지털 활용 격차를 설명하는 주된 요인임을 보여준다. \citet{Hargittai2010}는 인터넷 접근 이후의 격차, 즉 활용 방식의 차이가 실질적 불평등을 결정한다고 제시하였다. 사회경제적 지위에 따라 디지털 활용의 격차가 지속될 수 있음을 시사하는 결과이다. \citet{ChipevaEtAl2018}은 다양한 국가에서 실용성이 기술 사용의 가장 강력한 동기임을 보이는 한편, 개인의 내적 특성이 디지털 기기 사용에 유의미한 영향을 미친다는 점을 제시하였다. \citet{BorgesEtAl2020}은 앙골라 사례를 통해 기기 보급 이전에 습관과 ICT 기술 숙련도가 격차의 핵심 요인임을 보고하였다. \citet{Hilbert2011}의 라틴아메리카 18개국 분석은 남녀 간 디지털 격차가 존재하며, 국가 수준의 젠더 평등 지수가 높을수록 인터넷 사용이 활발함을 보였다.

한편, 개인 특성을 통제한 이후에도 지역 환경이 디지털 활용에 대해 추가적 설명력을 갖는다는 증거도 축적되어 있다. \citet{MossbergerEtAl2006}은 미국을 대상으로 HLM을 적용하여, 인터넷 접근 격차가 인종 자체보다 고빈곤 지역 환경에 기인함을 보였다. 이는 지역의 구조적 조건이 개인의 결과에 독립적 영향을 미칠 수 있음을 시사한다. \citet{CrangEtAl2006}은 영국 뉴캐슬의 질적 사례 비교를 통해, 지역 환경이 디지털 네트워크의 사용 방식 자체를 구조화함을 제시하였다. \citet{HongEtAl2017}의 중국 45세 이상 성인 분석은, 인터넷 이용자가 고소득, 고학력, 도시 거주 집단에 집중되어 있음을 보여, 지역 수준의 인프라와 기회 구조가 이용 가능성을 선별할 수 있음을 시사하였다.

\subsubsection{정책적 구분: 사람 기반 접근, 장소 기반 접근, 장소민감적 접근}
\label{subsubsec:2-1-2}

앞 절에서 검토한 두 유형의 요인은 공공정책 설계에서 사람 기반 접근(people-based approach)과 장소 기반 접근(place-based approach)의 구분에 대응한다. 사람 기반 정책은 개인의 역량, 소득, 교육 등 인적 특성을 직접 지원함으로써 결과를 개선하고자 한다. 장소 기반 정책은 특정 지역의 인프라, 제도, 환경 조건을 개선함으로써 그 지역 거주자의 결과를 간접적으로 향상시키고자 한다~\citep{NeumarkSimpson2015}. 두 접근은 상호 배타적이지 않으나, 정책 자원의 배분과 개입의 경로가 상이하므로 정책 수행의 관점에서 분리하여 설계될 필요가 있다.

그러나 두 접근 사이의 관계는 단순한 병렬이 아니다. \citet{Hindman2000}의 미국 시계열 분석은 소득, 연령, 교육이 지리보다 강력한 예측변수임을 보이면서도, 시간 경과에 따라 격차가 고착화되는 패턴을 보고하였다. 이는 개인 특성과 지역 구조가 시간 속에서 상호작용할 수 있음을 의미한다. \citet{vanDijk2020}는 동기 결여에서 물질적 접근, 기술 능력, 실질적 활용으로 이어지는 순차적 모형을 제시하여, 개인 요인과 구조 요인이 누적적으로 작용함을 제안하였다.

이러한 맥락에서, 최근 정책 및 학술 문헌은 사람 기반 또는 장소 기반의 양자택일이 아닌, 장소의 맥락 안에서 사람의 경험을 이해하는 장소민감적(place-sensitive) 접근으로 이행하고 있다. 생활 경험(lived experience)과 장소 기반 통찰(place-based insight)을 정량적 증거와 동등한 지식으로 인정하려는 시도가 그 예이다~\citep{OECD2025}. 이 접근은 사람 기반 요인과 장소 기반 요인을 분석적으로 분리하되, 두 요인이 특정 장소에서 어떻게 결합하는지를 함께 고려할 것을 요구한다.

\subsubsection{메타이론적 틀: 비판적 실재론의 층화된 존재론}
\label{subsubsec:2-1-3}

본 프로젝트는 사람 기반 요인과 장소 기반 요인의 분리를 분석 설계의 기본 원칙으로 채택하며, 그 설계의 메타이론적 기반을 비판적 실재론의 층화된 존재론에 둔다. \citet{Bhaskar1975}에 따르면 세계는 세 층위로 구성된다. 실재적 영역은 관찰 여부와 무관하게 존재하는 구조, 메커니즘, 인과적 권능을 포함한다. 실제적 영역은 이러한 구조와 메커니즘이 특정 조건에서 작동하여 사건을 발생시키는 영역이다. 경험적 영역은 발생한 사건 가운데 관찰 가능한 일부를 가리킨다. 실제적 영역의 모든 사건이 경험되는 것은 아니며, 관찰된 현상이 실재적 영역의 전부를 나타내는 것도 아니다.

이 존재론적 구분은 역행추론(retroduction)의 논리를 수반한다~\citep{Bhaskar1975}. 역행추론은 관찰된 경험적 사건으로부터, 그 사건을 가능하게 한 구조와 메커니즘을 역으로 추론하는 방법이다. 귀납이나 연역과 구분되며, \enquote{이 현상이 가능하려면 어떤 구조가 존재해야 하는가}를 묻는다.

\subsection{Seoul}
\label{subsec:2-2}

\subsubsection{도시로서의 서울}
\label{subsubsec:2-2-1}

대한민국의 수도인 서울은 국가 전체 국토 면적의 0.6\%에 인구의 20\%가 거주하는 거대한 밀집 도시이다. 높은 밀집도와 기술력의 집중에 따라, 서울은 대한민국 디지털 전환의 최전선에 존재해왔다. 그러나 서울의 내부를 면밀히 들여다보면 각 25개 자치구의 인구 구조와 경제적 자생력에 따른 유의미한 격차가 존재한다. 본 절에서는 서울 내 각 자치구별 특징을 조사하고, 서울의 인구 지형의 특징에 대해서 논의하고자 한다.

최근 서울의 인구 지형은 급격한 고령화와 1인 가구의 팽창이라는 두 가지 축으로 재편되고 있다. 강북구(26.03\%)와 도봉구(25.85\%) 등 서울 북부 외곽 지역은 고령 인구 비율이 이미 25\%를 상회하며 초고령사회적 특성을 띠는 반면, 강남권과 주요 대학가 주변은 상대적으로 젊은 인구 구조를 유지하고 있어 지역 간 연령 분포의 불균형이 뚜렷하다. 이러한 인구 통계적 차이는 각 구의 평균 소득 수준 및 경제적 활력과 결합하여, 본 연구 \ref{subsec:3-1}절에서 확인된 가격 부담(Affordability)이나 디지털 기기 보유(Devices) 역량의 차이를 발생시키는 구조적 배경이 된다.

\fig{image4.png}{서울 25개 자치구 인구\textperiodcentered{}연령 구조}{fig:2-pop}

또한, 서울의 정책 시행 체계는 서울특별시라는 광역 자치단체와 25개 기초 자치단체 사이의 유기적인 협력을 기반으로 작동한다. 정책의 큰 줄기는 서울시 본청에서 수립하는 하향식(Top-down) 방식을 따르는데, 이는 디지털 포용을 위한 광역 단위의 인프라 구축이나 표준 가이드라인 설정에 핵심적인 역할을 한다. 동시에 서울시는 시민참여예산제와 주민자치회 등을 통해 지역 주민이 실생활에서 느끼는 디지털 불편 사항을 정책에 직접 반영하는 상향식(Bottom-up) 의사결정 구조도 함께 운용하고 있다. 자치구는 이러한 광역의 계획을 바탕으로 지역 내 의회 심의와 조례 제정을 거쳐 예산을 확정하며, 최종적으로 동 주민센터를 통해 시민들에게 밀착형 서비스를 전달한다. 결국 서울의 디지털 격차 해소 정책은 중앙의 일괄적인 집행이 아닌, 각 자치구의 행정 역량과 시민 참여가 결합하여 구체화되는 다층적인 과정을 거쳐 실현된다고 볼 수 있다.

\subsubsection{연결성 격차 현황과 관련 법령 및 정책 현황}
\label{subsubsec:2-2-2}

서울의 디지털 연결성 실태는 자치구별 무선국 인프라 밀도의 편차를 통해 민간 주도 인프라 구축의 한계를 여실히 드러낸다. 전파누리의 2020--2025년 무선국 통계에 따르면, 서울의 통신 품질을 결정짓는 핵심 인프라인 '기지국'과 '이동중계국'은 수익성이 높고 트래픽이 집중되는 특정 지역에 압도적으로 쏠려 있다. 강남구(기지국 19,190개, 이동중계국 2,711개)와 송파구(기지국 14,751개) 등 동남권 지역은 촘촘한 망을 보유한 반면, 도봉구(기지국 4,008개)나 강북구(기지국 4,700개) 등 북부 외곽 지역은 기지국 수에서 최대 4.8배의 격차를 보인다. 특히 소규모 사업장과 사물인터넷(IoT) 연결에 활용되는 '간이무선국'의 경우 강남구(24,236개)가 도봉구(1,224개)보다 약 20배나 많아, 자치구 간 디지털 생태계의 부익부 빈익빈 현상이 물리적 인프라 수준에서 이미 고착화되고 있음을 시사한다.

\fig{image5.png}{자치구별 무선국 인프라 격차 (전파누리 2020--2025)}{fig:2-radio}

이러한 민간 인프라의 지역적 불균형을 교정하고 시민의 통신기본권을 보장하기 위해, 서울시는 행정적 대안으로서 직접 '기간통신사업자'로 등판하는 정책적 결단을 내렸다. 2025년 전기통신사업법 개정에 따라 지자체 최초로 통신사업자 지위를 확보한 서울시는, 민간 통신사가 수익성 문제로 투자를 미진하게 하는 소외 지역에 5,000km 이상의 '자가 광케이블망'을 직접 연결하여 공공 와이파이와 IoT 인프라를 확충하고 있다. 현재 서울시는 전국 최대 규모인 3만4천여 대의 공공와이파이를 운영하고 있으며, 이는 거리, 공원, 시장, 복지시설, 시내버스 등에 폭넓게 설치돼 연간 약 2천억 원의 시민 통신비 절감 효과를 거두고 있다. 특히 시내버스 7천여 대에 설치된 이동형 와이파이는 시민들이 일상에서 가장 체감하는 디지털 인프라로 자리잡고 있다. 서울시는 공공와이파이 확산을 통해 정보 접근의 격차를 줄이고, 매년 수천억 원에 달하는 가계 통신비 부담을 덜겠다는 계획이며, 또한 민간 임대망 의존도를 낮춰 확보된 재원을 시민 편익 사업에 재투자함으로써, 지속 가능한 선순환 구조를 마련하고자 하고 있다.

이러한 인프라 격차를 해소하기 위한 정책적 대응은 물리적 망 확충을 넘어, 시민의 실질적인 비용 부담을 낮추는 '경제적 접근성(Affordability)' 강화에 집중되고 있다. 정부와 지자체는 이용자가 자신의 통신 소비 패턴을 파악하고 합리적인 선택을 할 수 있도록 돕는 '스마트초이스(Smart Choice)' 플랫폼을 운영하고 있다. 스마트초이스는 단말기 지원금 조회, 맞춤형 요금제 추천, 미환급액 조회 등 통합 서비스를 제공하여 정보 비대칭성을 해소하고 가계 통신비 절감을 유도한다.

특히 최근 단행된 5G 요금제 개편은 기존 고가 요금제 중심의 시장 구조를 중저가 중심으로 전환했다는 점에서 큰 의의를 갖는다. 과거 4\,$\sim$\,5만 원대에 머물던 5G 최저 요금 구간을 3만 원대로 낮추었으며, 실제 데이터 사용량에 부합하도록 30\,$\sim$\,100GB 사이의 중간 요금제 구간을 촘촘하게 신설하였다. 이러한 개편은 이용자가 쓰지도 않는 데이터에 높은 비용을 지불하던 구조적 모순을 해결하고, 2024년 기준 600만 명 이상의 가입자가 중저가 요금제로 전환하는 실질적인 정책 성과로 이어졌다.

취약계층 및 특정 연령대를 위한 맞춤형 전용 요금제 역시 대폭 강화되었다. 대한민국 대표 이동통신사 3사에서는 청년층(만 34세 이하)을 위해, 데이터 수요가 많은 특성을 고려하여 일반 요금제와 동일한 가격에 데이터 제공량을 최대 2배까지 확대하고, 가입 연령 상한을 기존 만 29세에서 34세로 상향하여 수혜 대상을 넓혔다. 또한 LG U+에서는 고령층(만 65세 이상)을 위해, 상대적으로 데이터 사용량이 적고 가격 민감도가 높은 점을 반영하여, 일반 요금제 대비 최대 20\% 저렴한 전용 요금제를 연령대별(65세\textperiodcentered{}70세\textperiodcentered{}80세 등)로 세분화하여 출시하였다. 나아가 온라인 전용 요금제를 통해 비대면 가입 활성화 추세에 맞춰 대리점 방문 없이 온라인 채널로 가입 시 일반 요금제 대비 약 30\% 저렴하면서 약정이 없는 무약정 요금제를 도입하여 선택권을 넓혔다.

결과적으로 서울시는 스마트초이스를 통한 정보의 투명성, 통신 3사의 합리적 요금 체계, 그리고 서울시의 직접적인 공공 인프라라는 삼각 편대를 통해 모든 시민의 보편적 디지털 기본권을 보장하는 다각적인 포용 정책을 완성해 나가고 있다.

이에 더불어, 2026년 서울시 행정 서비스의 가장 큰 변화는 모바일 플랫폼의 단일화와 고도화이다. 서울시는 기존에 분절되어 운영되던 '서울시민카드'와 '서울지갑' 서비스를 2026년 3월 31일부로 종료하고, 통합 모바일 플랫폼인 '서울온(Seoul ON)'으로 기능을 일원화한다. 이는 시민들이 도서관, 체육시설 등 공공시설 이용부터 장애인\textperiodcentered{}국가유공자 자격 확인까지 하나의 앱에서 해결할 수 있게 함으로써, 서비스 파편화로 인한 디지털 행정의 진입장벽을 낮추고한 것이다. 특히 '서울온'은 공공 마이데이터와 연계하여 별도의 서류 제출 없이도 취약계층 자격을 실시간 확인하는 기능을 강화하여 행정 효율성을 높이는 기능을 한다.

나아가 서울시는 인프라 확충과 서비스 통합이 시민의 실질적인 이용 역량 강화로 이어질 수 있도록 'AI디지털배움터' 중심의 통합 포용 정책을 추진하고 있다. 서울 내 10개 거점을 포함하여, 전국 단위의 디지털 배움터 거점을 활용해 고령층 및 디지털 취약계층을 대상으로 한 대면 교육을 상시 운영하고 있으며, 단순한 인터넷 활용 교육을 넘어 실생활 중심의 AI 체험 교육과 디지털 역기능 예방 교육을 통합하여 제공한다. 특히 거동이 불편한 중증 장애인 등을 위해 '찾아가는 1:1 방문 교육'을 확대하고, 온\textperiodcentered{}오프라인이 연계된 수준별 맞춤형 교육 체계를 구축함으로써 교육 사각지대를 해소하고 있다. 이는 하드웨어(인프라)와 소프트웨어(행정 서비스)가 갖춰지더라도 이를 활용할 인적 역량이 부족하면 격차가 해소되지 않는다는 판단에 따른 것이다. 이처럼 서울시는 법령에 기반한 인프라 안정성 확보, '서울온'을 통한 서비스 전달 체계의 혁신, 그리고 교육을 통한 이용자 역량 제고라는 삼각 편대를 통해 보편적이고 의미 있는 연결성(UMC)을 실질적으로 구현해 나가고 있다.
